\section{江南供给与边区成本}

1132年临安逐渐成为行在、1138年确立为事实都城，1150年会子管理调整，1161年金军渡江受挫：这些事件显示江南城市、纸币和军事供给彼此相连。它们不证明“江南财富”可以无成本地供养全部战线。临安的漕运和市场、两淮的边防、四川的道路、福建与两广的海域各有不同的征收、运输与风险。

契约、讼牍、港口与窑址材料能呈现局部的土地、借贷、工艺和贸易，却必须先辨年代、原件或抄件、当事人和出土地点。一艘沉船或一批外销瓷不等于国家直接经营海贸；一纸田契也不等于普遍土地关系。把中央会要、地方材料与考古并置，才能理解南宋财政如何依赖区域差异，而不是抹去它们。
